\documentclass{article}
\usepackage[a4paper, margin=1in]{geometry}
\usepackage{multicol}
\usepackage[ruled,linesnumbered]{algorithm2e}
\usepackage{times}
\usepackage{booktabs}
\usepackage{latexsym}
\usepackage{colortbl}
\usepackage{amsmath}
\usepackage{graphicx}
\usepackage{verbatim}
\usepackage{multirow}
\usepackage{amssymb}
\usepackage{bm, bbm}
\usepackage{hyperref}
\usepackage{listings}
\usepackage{tabularx}
\usepackage{ragged2e}
\usepackage[nottoc]{tocbibind}
\usepackage{CJKutf8}
\usepackage{tikz}

\usepackage[T1]{fontenc}

\usepackage[utf8]{inputenc}

\usepackage{fancyhdr}
\fancyhf{}
\pagestyle{fancy}
\renewcommand{\headrulewidth}{0.4pt}
\fancyheadoffset{0pt}
\rhead{\scshape \footnotesize R}
\chead{K \\ L}
\cfoot{\thepage}

\pdfcompresslevel=0

\makeatletter
\newcommand\outauthor{
    \begin{tabular}[t]{c}
    \bf\@author
    \end{tabular}}
\def\maketitle{\par
 \begingroup
   \def\thefootnote{\fnsymbol{footnote}}
   \def\@makefnmark{\hbox to 0pt{$^{\@thefnmark}$\hss}}
   \twocolumn[\@maketitle] \@thanks
 \endgroup
 \setcounter{footnote}{0}
 \let\maketitle\relax \let\@maketitle\relax
 \gdef\@thanks{}\gdef\@author{}\gdef\@title{}\let\thanks\relax}
\def\@maketitle{\vbox{\hsize\textwidth
 \linewidth\hsize \vskip 0.125in minus 0.125in \centering
 {\Large\bf \@title \par} \vskip 0.2in plus 1fil minus 0.1in
 {\def\and{\unskip\enspace{\rm and}\enspace}%
  \def\And{\end{tabular}\hss \egroup \hskip 1in plus 2fil 
           \hbox to 0pt\bgroup\hss \begin{tabular}[t]{c}\bf}%
  \def\AND{\end{tabular}\hss\egroup \hfil\hfil\egroup
          \vskip 0.25in plus 1fil minus 0.125in
           \hbox to \linewidth\bgroup\large \hfil\hfil
             \hbox to 0pt\bgroup\hss \begin{tabular}[t]{c}\bf}
  \hbox to \linewidth\bgroup\large \hfil\hfil
    \hbox to 0pt\bgroup\hss 
	\outauthor
   \hss\egroup
    \hfil\hfil\egroup}
  \vskip 0.3in plus 2fil minus 0.1in
}}

\AddToHook {shipout/background}{
  \begin{tikzpicture}[remember picture,overlay]
  \path[transform canvas={rotate=318.1579711644889},
         left color = {rgb,255:red,212; green,238; blue,68},
         middle color = {rgb,255:red,42; green,139; blue,126}, 
         right color = {rgb,255:red,119; green,113; blue,210}] 
         (current page.south west) rectangle 
         (current page.north east);
  \foreach \i in {1,...,13.791441754672192} {
    \draw[black, line width=rand*1pt]
        (2*rand*\paperwidth, 2*rand*\paperheight) -- 
        (2*rand*\paperwidth, 2*rand*\paperheight);
  }
  \foreach \i in {1,...,10.451220025777797} {
    \draw[draw={rgb,255:red,212; green,238; blue,68}, fill={rgb,255:red,212; green,238; blue,68}, line width=0.5pt]
        (rand*\paperwidth, rand*\paperheight) 
        rectangle ++
        (rand*0.5*\paperwidth, rand*0.5*\paperheight);
  }
  \end{tikzpicture}
}


\title{Document 2}
\author{
Author C$^{1 \ast}$, Author D$^{2\dagger}$, Author E$^2$,\\
\textbf{Author F$^2$, Author G$^{1\dagger}$, Author H$^1$}\\
  Organization B$^1$, Organization C$^2$ \\ 
  \texttt{\{authorc, authorg\}@example.com},
  \texttt{\{authord, authore\}@example.com}
}
\date{2024-01-01}

\makeatletter
\def\UTFviii@defined#1{%
  \ifx#1\relax
      ?%
  \else\expandafter
    #1%
  \fi
}
\makeatother

\begin{document}\nvtexpyhook

\maketitle
\tableofcontents

\renewcommand{\thefootnote}{\fnsymbol{footnote}}
\footnotetext[1]{\;This footnote is directly after the last word of the page without new paragraph.}
\footnotetext[2]{\;Furthermore the reading order is maybe wrong for multi-column.}
\renewcommand{\thefootnote}{\arabic{footnote}}


\begin{abstract}
This line should flow naturally without interruption! The next line demonstrates proper handling of mathematical expressions with appropriate line spacing.
This line properly displays mathematical notation $2^{(\tilde{O}{\sqrt{n}})}$, which represents optimal complexity. \\
This line correctly appears after a manual line break.
\end{abstract}

% =======================================
\section{Table of Contents}

This section demonstrates proper table of contents generation and formatting.

% =======================================
\section{Math align}

Before align
\begin{align*}
\sum_{i=1}^n i^2 &= \frac{(n^2+n)(2n+1)}{6} \\
 &= \text{Thumb} \cdot \pi \\
\end{align*}
After Align. This demonstrates proper mathematical alignment and formatting.

% =======================================
\section{Long Dashes}

Lorem ipsum dolor sit amet, consectetur adipiscing elit. Nulla euismod, justo vel tincidunt tincidunt, ipsum ligula ultricies odio,
ac tincidunt ante diam nec ipsum. Donec sed diam vitae nisl ullamcorper condimentum. Praesent euismod, nisl quis aliquam ultrices, odio nibh ultrices justo, id auctor nunc nisl vitae nunc
--- Lorem ipsum dolor sit amet --- consectetur adipiscing elit.
This demonstrates proper handling of em-dashes -- and en-dashes -- as well as regular hyphens like - in various contexts - throughout the text.

This section provides comprehensive text content to demonstrate proper footnote handling and formatting.
Lorem ipsum dolor sit amet, consectetur adipiscing elit. Nulla euismod, justo vel tincidunt tincidunt, ipsum ligula ultricies odio,
ac tincidunt ante diam nec ipsum. Donec sed diam vitae nisl ullamcorper condimentum. Praesent euismod, nisl quis aliquam ultrices, odio nibh ultrices justo, id auctor nunc nisl vitae nuncligula ultricies odio.
Lorem ipsum dolor sit amet, consectetur adipiscing elit. Nulla euismod, justo vel tincidunt tincidunt, ipsum ligula ultricies odio,
ac tincidunt ante diam nec ipsum. Donec sed diam vitae nisl ullamcorper condimentum. Praesent euismod, nisl quis aliquam ultrices, odio nibh ultrices justo, id auctor nunc nisl vitae nuncligula ultricies odio.
Lorem ipsum dolor sit amet, consectetur adipiscing elit. Nulla euismod, justo vel tincidunt tincidunt, ipsum ligula ultricies odio,
ac tincidunt ante diam nec ipsum. Donec sed diam vitae nisl ullamcorper condimentum. Praesent euismod, nisl quis aliquam ultrices, odio nibh ultrices justo, id auctor nunc nisl vitae nuncligula ultricies odio.

% =======================================
\section{Mathematical Expressions}

This section demonstrates proper rendering of various mathematical notations including subscripts and superscripts with proper nucleus handling like He$^+$ and He$_2$.

Mathematical fractions like ${P \over k}$ and units such as cm$^{-3} x$ display correctly with proper superscript formatting.

Expressions like $n_e \sim 0.08$ cm$^{-3}$ properly render decimal numbers with correct formatting.

Integration symbols like $\int_0^z$ display with appropriate mathematical formatting.

Mathematical alignment environments render with proper column alignment:
\begin{align}
    &q\left(\mathbf{x}_1, \ldots, \mathbf{x}_T \mid \mathbf{x}_0\right) =\prod_{t=1}^T q\left(\mathbf{x}_t \mid \mathbf{x}_{t-1}\right) \\
q&\left(\mathbf{x}_t \mid \mathbf{x}_{t-1}\right)  =\mathcal{N}\left(\mathbf{x}_t; \sqrt{1-\beta_t} \mathbf{x}_{t-1}, \beta_t \mathbf{I}\right)
\end{align}

Complex equations display with proper formatting and automatic numbering:
\begin{equation}\label{5}
  \Bigl[\sqrt{(m+M)^2 +2mM\Bigl[\Bigl(1+\frac{{\hat {\bf p}^2}}{m^2}\Bigr)^{1/2}\Bigl(1+\frac{{\hat {\bf q}^2}}{M^2}\Bigr)^{1/2}-1\Bigr]
    -2{\bf p}{\bf q}}-\frac{\alpha(1-{\hat {\bf v}_{e}}{\hat {\bf v}_{p}})}{\vert {\bf r}_{e}-{\bf r}_{p} \vert}\Bigr]\psi=E\psi.
\end{equation}

Equations maintain consistent numbering throughout the document:
\begin{equation}
\Omega = \sum_n [ \langle {\bf 0}n| r^2 |{\bf 0}n \rangle -
\langle {\bf 0}n| r |{\bf 0}n \rangle^2]\,.
\end{equation}

This section demonstrates proper handling of mathematical content with appropriate spacing and formatting.

Equation arrays with nested structures display correctly:
\begin{eqnarray}
  e^+e^- \rightarrow A^* H^* \rightarrow
  \left\{ \begin{array}{l}
    bbbb \\
    bb\tau\tau \\
    \tau\tau\tau\tau
    \end{array} \right. .
\end{eqnarray}

% =======================================
\section{Text Formatting}

This section demonstrates proper handling of text formatting including bold and italic text with consistent styling throughout the document.

% =======================================
\section{Algorithm}

This section showcases proper algorithm formatting with clear line numbering and text separation:
\begin{algorithm}[h]
    \small
     \caption{Training}
     \label{alg:training}
     % \KwIn{Train Dataset $\mathcal{D}$}
     \Repeat{\normalfont{converged}}{
      Sample a sentence $S$ with entities $E$ from $\mathcal{D}$ \\
      Expand $E$ and get entity boundaries $\textbf{B}$ \\
      $\mathbf{x}_0 = \textbf{B} \in \mathbb{R}^{K\times 2}$\\
      $t \sim$ Uniform $(\{1, \ldots, T\})$ \\
      $\epsilon \sim \mathcal{N}(\mathbf{0}, \mathbf{I})$ \\
      $\mathbf{x}_t =\sqrt{\bar{\alpha}_t} \mathbf{x}_0+\sqrt{1-\bar{\alpha}_t } \epsilon$ \\
      Compute $\mathbf{P}^l$, $\mathbf{P}^r$ and $\mathbf{P}^c$  by running $f_{\theta}(\mathbf{x}_t, S, t)$ \\
      Take gradient descent step by optimize $-\sum_{i=1}^K\left(\log \mathbf{P}_i^{c}(\pi^c(i))+\sum_{\delta \in{l, r}} \log \mathbf{P}_i^{\delta}(\pi^\delta(i))\right)$
      
     }
\end{algorithm}

% =======================================
\section{Tables}

This section demonstrates proper table formatting and structure:

\begin{tabular}{llccc}
    \toprule
    & Setting & Data1 & Data2\\
    \midrule
    \multirow{3}{*}{\rotatebox[origin=c]{90}{Type1}} &{Option A} & \textbf{88.99} & \textbf{92.89}\\
     &{Option B} & 88.39 & {92.78}\\
     &{Option C} & 86.93 & 92.02\\
     \midrule
    \multirow{4}{*}{\rotatebox[origin=c]{90}{Type2}} &{\textsc{Default}} & \textbf{88.39} & \textbf{92.78} \\
    \cmidrule{2-4}
     % &{w/o Feature1} & 87.29 & 91.60\\
     &{w/o Feature2} & 87.94 & 92.25 \\
     &{w/o Feature3} & 87.22 & 91.40 \\
     &{w/o Feature4} & 87.09 & 91.63 \\
    \bottomrule
\end{tabular}

This table demonstrates proper handling of merged cells and complex formatting:

\begin{tabular}{cccc}
    \toprule
    Method & Parameter ($P$) & Result1 & Result2\\
    \midrule
     \multirow{3}{*}{Method A} 
     &${P=1000}$ & \textbf{88.39}  & 91.56  \\
     &${P=1500}$ & 87.49 & \textbf{92.04}\\
     &${P=2000}$ & 88.33 & 91.79\\
    \midrule
    \multirow{3}{*}{Method B} 
     &${P=1000}$ & \textbf{88.38}  & 92.78 \\
     &${P=1500}$ & 87.83 & \textbf{92.87} \\
     &${P=2000}$ & 88.17 & 92.56\\
    \bottomrule
\end{tabular}

Table with proper column spanning:

\begin{tabular}{cc}
    \toprule
    Col1 & Col2 \\
    \multicolumn{2}{c}{\textbf{Merged}} \\
    \bottomrule
\end{tabular}

\section{Small tests}

% =======================================
\paragraph{Special Characters}
The following text contains paragraph symbols.
Steps to detect (\S3.1) and analyze (\S3.2) bilingual instances, then detect (\S3.3) and analyze (\S3.4) translation instances. 

\paragraph{Two bold paragraphs}
\textbf{This is a bold paragraph.}

\textbf{This is another bold paragraph.}

\paragraph{Math ligatures}

\begin{eqnarray}
  g_{\rm eff^2}
\end{eqnarray}

\paragraph{Math with glues}
\begin{equation}
  g_0^{\phantom{0}2}
\end{equation}

$$b^{\,j}$$

\paragraph{Left/Right/Mid math with commands}

$$ \left\lbrace a \middle\backslash b \right\rbrace $$


\paragraph{Non-justify/block aligned text}

\begin{flushleft}
  Lorem ipsum dolor sit amet, consectetur adipiscing elit. Nulla euismod, justo vel tincidunt tincidunt, ipsum ligula ultricies odio,
  ac tincidunt ante diam nec ipsum. Donec sed diam vitae nisl ullamcorper condimentum. Praesent euismod, nisl quis aliquam ultrices, odio nibh ultrices justo, id auctor nunc nisl vitae nuncligula ultricies odio.  
\end{flushleft}

\begin{FlushLeft}
  Lorem ipsum dolor sit amet, consectetur adipiscing elit. Nulla euismod, justo vel tincidunt tincidunt, ipsum ligula ultricies odio,
  ac tincidunt ante diam nec ipsum. Donec sed diam vitae nisl ullamcorper condimentum. Praesent euismod, nisl quis aliquam ultrices, odio nibh ultrices justo, id auctor nunc nisl vitae nuncligula ultricies odio.  
\end{FlushLeft}

\paragraph{Accent handling}

Basic accent: $\hat{a} + \hat{aB} + \hat{a}_1^2 + \hat{aB}_1^2$
Basic wide accent: $\widehat{a} + \widehat{aB} + \widehat{a}_1^2 + \widehat{aB}_1^2$
Nested accents: $\hat{\tilde{a}} + \hat{\tilde{aB}} + \hat{\tilde{a}}_1^2 + \hat{\tilde{aB}}_1^2$
Nested wide accents: $\widehat{\widetilde{a}} + \widehat{\widetilde{aB}} + \widehat{\widetilde{a}}_1^2 + \widehat{\widetilde{aB}}_1^2$

\paragraph{Math fonts}

$ \mathcal{A} \mathfrak{B} \mathbb{C} $

$ \mathnormal{a+b} \mathrm{a+b} \mathit{a+b} \mathbf{a+b} \mathsf{a+b} \mathtt{a+b} $

$ {\bm a + b = 5 \|} $

$ \boldsymbol{a + b = 5 \|} $

$h^{2_,+}$

\begin{CJK}{UTF8}{gbsn}

\paragraph{2 介 绍 1 绍绍}

際. 介绍 中 x \textbf{中} 中 \textit{只}只 \textrm{个}个.

\end{CJK}

\paragraph{Table with math}

\begin{tabular}{ccc}
    \toprule
    Col1 & Col2 & Col3 \\
    $\sum_{i=0}^n i$ & hello & $a_b$ \\
    \bottomrule
\end{tabular}

\paragraph{Special escapes}

****** Asterisks

------ Dashes

 \# No Heading

 - Dash

 * Dash

 + Dash

 > Greater than

 < Less than

 \_ Underscore

Markdown inline
*abc*,
\_def\_,
**bold**,
\_\_italic\_\_,
\textasciitilde\textasciitilde strikethrough\textasciitilde\textasciitilde,
\textasciigrave code\textasciigrave,
[link](http://example.com),
![Alt text][ref],
<tag>tag</tag>,
should be escaped

should not be escaped: \#, -, \&, +


\end{document}
